\documentclass[11pt,a4paper]{article}
\usepackage[margin=25mm]{geometry}
\usepackage[T1]{fontenc}
\usepackage{lmodern}
\usepackage{amsmath,amssymb}
\usepackage{booktabs}
\usepackage{microtype}
\usepackage[colorlinks=true,linkcolor=blue,urlcolor=blue]{hyperref}

\newcommand{\E}{\mathbb{E}}
\newcommand{\sg}{\operatorname{stopgrad}}
\newcommand{\Var}{\operatorname{Var}}
\newcommand{\as}{\xrightarrow{\mathrm{a.s.}}}
\setlength{\parindent}{0pt}
\setlength{\parskip}{6pt}
\setlength{\emergencystretch}{2em}
\allowdisplaybreaks[1]
\hypersetup{
  pdftitle={Deriving difficulty reweightings: CS336 Assignment 5},
  pdfsubject={Section 5.3: Deriving difficulty reweightings},
  pdfauthor={Study notes}
}

\begin{document}

\begin{center}
  {\LARGE\bfseries Deriving difficulty reweightings}\par
  \vspace{4pt}
  {\large CS336 Assignment 5, Section 5.3}\par
  \texttt{derive\_difficulty\_reweightings}
\end{center}

\textbf{Source and scope.} These notes follow the Spring 2026 handout,
pp.~28--29, especially its equations (39)--(43). Throughout, rewards are binary,
the loss normalizer is the constant $Z=G$, and we take $G\to\infty$.
Thus part (b) concerns GRPO with \emph{constant} loss normalization.

Write $\eta_\theta(x)=\E_{y\sim\pi_\theta(\cdot\mid x)}[r(y\mid x)]$.
The answers are the following prompt weights, evaluated at the current policy
and held fixed when taking a gradient:
\begin{center}
\renewcommand{\arraystretch}{1.6}
\begin{tabular}{lll}
\toprule
Part & Method & $w\bigl(x,\sg(\pi_\theta)\bigr)$ \\
\midrule
(a) & Dr. GRPO & $1$ \\
(b) & GRPO, constant loss normalization & $\displaystyle \frac{1}{\sqrt{\eta_\theta(x)(1-\eta_\theta(x))}}$ \\
(c) & MaxRL & $\displaystyle \frac{1}{\eta_\theta(x)}$ \\
\bottomrule
\end{tabular}
\end{center}
For (b), assume $0<\eta_\theta(x)<1$; for (c), assume $\eta_\theta(x)>0$.
The last page discusses zero denominators and the finite-group case.

\textbf{Notation.} Fix a prompt $x$ and the current parameters $\theta_0$.
Sample $Y_1,\ldots,Y_G$ independently from $\pi_{\theta_0}(\cdot\mid x)$, and set
\begin{equation}
  R_j=r(Y_j\mid x),\qquad
  S_j=\left.\nabla_\theta\log\pi_\theta(Y_j\mid x)\right|_{\theta=\theta_0},\qquad
  \mu_G=\frac{1}{G}\sum_{j=1}^{G}R_j.
\end{equation}
Here $S_j$ is a vector, and $\eta=\eta_{\theta_0}(x)$ is the probability of a
correct response. Higher $\eta$ means an easier prompt for this policy.
Assume that the prompt distribution and reward function do not depend on
$\theta$, and that differentiation and expectation can be interchanged.

\textbf{Two identities used in every part.} The score has conditional mean zero:
\begin{align}
  \E[S_j\mid x]
  &=\left.\sum_y \pi_\theta(y\mid x)
                  \nabla_\theta\log\pi_\theta(y\mid x)\right|_{\theta_0} \notag\\
  &=\left.\nabla_\theta\sum_y\pi_\theta(y\mid x)\right|_{\theta_0}
   =\nabla_\theta 1=0. \label{eq:score-zero}
\end{align}
Multiplying by reward gives the policy-gradient identity:
\begin{align}
  \E[R_jS_j\mid x]
  &=\left.\sum_y r(y\mid x)\nabla_\theta\pi_\theta(y\mid x)\right|_{\theta_0}
   =\left.\nabla_\theta\eta_\theta(x)\right|_{\theta_0}.
   \label{eq:reward-score}
\end{align}
Consequently, subtracting the population mean does not change this expectation:
\begin{equation}
  \E[(R_j-\eta)S_j\mid x]
  =\E[R_jS_j\mid x]-\eta\E[S_j\mid x]
  =\left.\nabla_\theta\eta_\theta(x)\right|_{\theta_0}.
  \label{eq:baseline}
\end{equation}

\newpage

\textbf{One calculation for all three normalizers.}
Let $d_G$ denote the group normalizer: $1$, the group standard deviation, or
$\mu_G$. At this fixed prompt, the estimator is
\begin{align}
  \widehat g_G(x)
  &=\frac{1}{G}\sum_{j=1}^{G}\frac{R_j-\mu_G}{d_G}S_j \notag\\
  &=\frac{1}{d_G}
    \left(\frac{1}{G}\sum_{j=1}^{G}R_jS_j
          -\mu_G\frac{1}{G}\sum_{j=1}^{G}S_j\right).
    \label{eq:sample-rewrite}
\end{align}
Suppose $d_G\as d(x)>0$. The law of large numbers, applied to each vector
coordinate, gives
\begin{equation}
  \mu_G\as\eta,\qquad
  \frac{1}{G}\sum_j S_j\as 0,\qquad
  \frac{1}{G}\sum_j R_jS_j
    \as\left.\nabla_\theta\eta_\theta(x)\right|_{\theta_0}.
\end{equation}
Substituting these limits into \eqref{eq:sample-rewrite},
\begin{equation}
  \boxed{\displaystyle
  \widehat g_G(x)\as
    \frac{1}{d(x)}
    \left.\nabla_\theta\eta_\theta(x)\right|_{\theta_0}.}
  \label{eq:master}
\end{equation}
The required moments must exist; averaging this limit over prompts also
requires the corresponding integrability/interchange conditions.

To match the handout's stop-gradient surrogate, freeze the reference policy:
\begin{equation}
  J_w(\theta;\theta_0)
  =\E_{x\sim\rho}\!\left[w(x,\pi_{\theta_0})\eta_\theta(x)\right].
  \label{eq:frozen-objective}
\end{equation}
Its gradient at the current policy is
\begin{equation}
  \left.\nabla_\theta J_w(\theta;\theta_0)\right|_{\theta_0}
  =\E_{x\sim\rho}\!\left[
       w(x,\pi_{\theta_0})
       \left.\nabla_\theta\eta_\theta(x)\right|_{\theta_0}\right].
  \label{eq:surrogate-gradient}
\end{equation}
Comparing \eqref{eq:master} and \eqref{eq:surrogate-gradient} gives $w=1/d$.

\textbf{(a) Dr. GRPO.} Here $d_G=1$, so $d(x)=1$ and
\begin{equation}
  \boxed{w_{\mathrm{Dr.\,GRPO}}(x,\pi_{\theta_0})=1.}
\end{equation}
The mean baseline cancels in expectation by \eqref{eq:baseline}. In the
infinite-group limit, the estimator is therefore the gradient of ordinary
expected reward, with the original prompt distribution $\rho$.

\textbf{(b) GRPO with constant loss normalization.} Since $R_j\in\{0,1\}$,
$R_j^2=R_j$. Hence
\begin{equation}
  \Var(R_j\mid x)=\E[R_j^2\mid x]-\E[R_j\mid x]^2
  =\eta-\eta^2=\eta(1-\eta).
\end{equation}
The group standard deviation converges to $\sqrt{\eta(1-\eta)}$. This is true
whether the sample variance uses denominator $G$ or $G-1$. Thus
\begin{equation}
  \boxed{\displaystyle
  w_{\mathrm{GRPO}}(x,\pi_{\theta_0})
    =\frac{1}{\sqrt{\eta_{\theta_0}(x)(1-\eta_{\theta_0}(x))}}.}
\end{equation}

\textbf{(c) MaxRL.} Here $d_G=\mu_G\as\eta$. Therefore
\begin{equation}
  \boxed{\displaystyle
  w_{\mathrm{MaxRL}}(x,\pi_{\theta_0})
    =\frac{1}{\eta_{\theta_0}(x)}.}
\end{equation}
Equivalently, its population gradient is
$\E_{x\sim\rho}[\nabla_\theta\log\eta_\theta(x)]$ at $\theta=\theta_0$,
where the log and expectation are well-defined.

\newpage

\textbf{Why \texttt{stopgrad} matters.}
In \eqref{eq:frozen-objective}, $\theta$ is differentiated and $\theta_0$ is held
fixed. This is precisely the meaning of the handout's
$w(x,\sg(\pi_\theta))$. The derivative does not include an extra
$\eta_\theta(x)\nabla_\theta w$ term.

MaxRL makes the distinction especially clear:
\begin{align}
  J_{\mathrm{MaxRL}}(\theta;\theta_0)
    &=\E_x\!\left[\frac{\eta_\theta(x)}{\eta_{\theta_0}(x)}\right],\\
  J_{\mathrm{MaxRL}}(\theta_0;\theta_0)&=1,\\
  \left.\nabla_\theta J_{\mathrm{MaxRL}}(\theta;\theta_0)\right|_{\theta_0}
    &=\E_x\!\left[\frac{\left.\nabla_\theta\eta_\theta(x)\right|_{\theta_0}}
                         {\eta_{\theta_0}(x)}\right].
\end{align}
Although the surrogate's value at the current parameters is $1$, its derivative
with respect to the first argument need not be zero. Simplifying the expression
to $\eta_\theta/\eta_\theta=1$ before differentiating would incorrectly let the
denominator vary. For a scalar illustration,
$f(p;p_0)=p/p_0$ satisfies $f(p_0;p_0)=1$, but
$\partial f/\partial p=1/p_0$.

\textbf{Connection to the finite-group baseline bias.}
At finite $G$, $\mu_G$ contains $R_j$, so it is not independent of $S_j$.
Using independence only for the other samples,
\begin{align}
  \E[\mu_GS_j\mid x]
  &=\frac{1}{G}\E[R_jS_j\mid x]
    +\frac{1}{G}\sum_{k\ne j}\E[R_k\mid x]\E[S_j\mid x] \notag\\
  &=\frac{1}{G}\left.\nabla_\theta\eta_\theta(x)\right|_{\theta_0}.
\end{align}
Therefore the finite-group Dr. GRPO estimator has expectation
\begin{equation}
  \E[\widehat g_G^{\mathrm{Dr.\,GRPO}}(x)\mid x]
  =\left(1-\frac{1}{G}\right)
    \left.\nabla_\theta\eta_\theta(x)\right|_{\theta_0}.
\end{equation}
The factor tends to $1$, explaining why it does not appear in part (a).
For (b) and (c), the denominator is also random and correlated with the
numerator at finite $G$. Replacing an expectation of a ratio by a ratio of
expectations is unjustified; the preceding derivation uses the large-group
limit instead.

\textbf{Reading the weights.}
Dr. GRPO adds no difficulty-dependent multiplier. GRPO's multiplier is
smallest at $\eta=1/2$ and grows toward both endpoints. MaxRL's multiplier
decreases as $\eta$ grows, emphasizing harder prompts relative to easier ones.
For example:
\begin{center}
\begin{tabular}{lrrr}
\toprule
Success probability $\eta$ & Dr. GRPO & GRPO & MaxRL \\
\midrule
$0.1$ & $1$ & $3.33$ & $10$ \\
$0.5$ & $1$ & $2$ & $2$ \\
$0.9$ & $1$ & $3.33$ & $1.11$ \\
\bottomrule
\end{tabular}
\end{center}
These numbers multiply each prompt's own gradient $\nabla_\theta\eta_\theta(x)$.
A larger weight alone does not imply a larger resulting gradient norm.

\textbf{Zero denominators and implementation.}
Without stabilization, GRPO is undefined at $\eta=0$ or $1$, and MaxRL is
undefined at $\eta=0$. Even when $0<\eta<1$, a finite group can contain identical
rewards; all centered advantages then vanish. Adding $\varepsilon>0$ to the
normalizer makes those groups contribute zero. For fixed $\varepsilon$, the
population weights become
\begin{equation}
  w_{\mathrm{GRPO},\varepsilon}(x)
    =\frac{1}{\sqrt{\eta(x)(1-\eta(x))}+\varepsilon},
  \qquad
  w_{\mathrm{MaxRL},\varepsilon}(x)
    =\frac{1}{\eta(x)+\varepsilon}.
\end{equation}
The boxed answers use the problem's unstabilized formulas. They also rely on
constant loss normalization: dividing each response by its own length
introduces a response-dependent factor, which this derivation does not remove.

\end{document}
